\section*{Capítulo \ref{ch:g11:reaction-energy} --- A energia das reações: combustão e energias de ligação}

\begin{solution}{exo:g11:reaction-energy:1}
A madeira queimando: \omterm{def:g11:reaction-energy:exothermic}{exotérmica}. O gelo derretendo: \omterm{def:g11:reaction-energy:exothermic}{endotérmica} (ele
tira calor da vizinhança, embora não seja uma \omterm{def:g7:chemical-reactions:reaction}{reação química}). Compressa
fria: \omterm{def:g11:reaction-energy:exothermic}{endotérmica}. Aquecedor de mãos: \omterm{def:g11:reaction-energy:exothermic}{exotérmica}. Produzir açúcar à luz
do Sol: \omterm{def:g11:reaction-energy:exothermic}{endotérmica} (a energia vem da luz).
\end{solution}

\begin{solution}{exo:g11:reaction-energy:2}
$Q = 3.0 \times 890.7 = \qty{2.7e3}{kJ}$, cerca de \qty{2.7}{MJ}.
\end{solution}

\begin{solution}{exo:g11:reaction-energy:3}
Os \omterm{def:g7:chemical-reactions:reactant}{reagentes}. A diferença é liberada para a vizinhança, principalmente
na forma de calor.
\end{solution}

\begin{solution}{exo:g11:reaction-energy:4}
$n = 100 / 46.0 = \qty{2.17}{mol}$; $Q = 2.17 \times 1367.6 =
\qty{2.97e3}{kJ}$, cerca de \qty{3.0}{MJ}.
\end{solution}

\begin{solution}{exo:g11:reaction-energy:5}
A energia necessária para romper um \omterm{def:g10:the-mole:amount}{mol} dessas ligações, com as
\omterm{def:g7:atoms-and-molecules:molecule}{moléculas} e os \omterm{def:g7:atoms-and-molecules:atom}{átomos} no estado gasoso. O par compartilhado mantém os
dois \omterm{def:g7:atoms-and-molecules:atom}{átomos} juntos: afastá-los é agir contra essa atração, o que exige
energia.
\end{solution}

\begin{solution}{exo:g11:reaction-energy:6}
Rompidas: $2 \times 436 + 498 = \qty{1370}{kJ}$; formadas:
$4 \times 464 = \qty{1856}{kJ}$; $E_r \approx 1370 - 1856 =
\qty{-486}{kJ}$: \omterm{def:g11:reaction-energy:exothermic}{exotérmica}.
\end{solution}

\begin{solution}{exo:g11:reaction-energy:7}
O etanol \ce{CH3-CH2-O-H} tem 5 \ce{C-H}, 1 \ce{C-C}, 1 \ce{C-O} e 1
\ce{O-H}. Rompidas: $5 \times 415 + 345 + 350 + 464 + 3 \times 498 =
\qty{4728}{kJ}$. Formadas: $4 \times 741 + 6 \times 464 = \qty{5748}{kJ}$.
$E_r \approx \qty{-1020}{kJ}$, cerca de um quarto menor, em valor
absoluto, que os \qty{1367.6}{kJ} medidos.
\end{solution}

\begin{solution}{exo:g11:reaction-energy:8}
Propano: $2219.2 / 0.0440 = \qty{5.04e4}{kJ/kg} = \qty{50.4}{MJ/kg}$.
Metano: $890.7 / 0.0160 = \qty{55.7}{MJ/kg}$. Os dois como no gráfico.
\end{solution}

\begin{solution}{exo:g11:reaction-energy:9}
$Q = 1500 \times 4.18 \times 85 = \qty{5.33e5}{J} = \qty{533}{kJ}$;
$n = 533 / 890.7 = \qty{0.598}{mol}$, isto é, $0.598 \times 16.0 =
\qty{9.6}{g}$ de metano; com metade da energia perdida, \qty{19}{g}.
\end{solution}

\begin{solution}{exo:g11:reaction-energy:10}
Hidrogênio > metano > propano > octano > etanol.
$48.0 / 142.9 = \qty{0.34}{kg}$ de hidrogênio.
\end{solution}

\begin{solution}{exo:g11:reaction-energy:11}
Rompidas: $946 + 3 \times 436 = \qty{2254}{kJ}$; formadas:
$6 \times 390 = \qty{2340}{kJ}$; $E_r \approx \qty{-86}{kJ}$: levemente
\omterm{def:g11:reaction-energy:exothermic}{exotérmica}.
\end{solution}

\begin{solution}{exo:g11:reaction-energy:12}
Água: $250 \times 4.18 \times 20 = \qty{2.09e4}{J} = \qty{20.9}{kJ}$.
Etanol: $1.20 / 46.0 = \qty{0.0261}{mol}$, que poderia liberar
$0.0261 \times 1367.6 = \qty{35.7}{kJ}$. Eficiência
$20.9 / 35.7 \approx \qty{59}{\percent}$. Perdas: calor levado pelo \omterm{def:g4:air-a-mixture-of-gases:air}{ar},
calor que aquece a lata e o suporte, \omterm{def:g8:combustion-and-fuels:complete}{combustão incompleta} (fuligem),
uma parte do etanol que \omterm{def:g3:separating-mixtures:evaporate}{evapora} sem queimar.
\end{solution}

\begin{solution}{exo:g11:reaction-energy:13}
\ce{2C8H18 + 25O2 -> 16CO2 + 18H2O}: 8 \omterm{def:g10:the-mole:amount}{mols} de dióxido de carbono,
\qty{352}{g}, por \qty{5470}{kJ}, logo $352 / 5.470 = \qty{64.4}{g}$ por
megajoule. \ce{C3H8 + 5O2 -> 3CO2 + 4H2O}: \qty{132}{g} por
\qty{2219.2}{kJ}, logo \qty{59.5}{g} por megajoule. O metano, com
\qty{49.4}{g}, é o que libera menos.
\end{solution}

\begin{solution}{exo:g11:reaction-energy:14}
A tabela dá \omterm{def:g11:reaction-energy:bond-energy}{energias de ligação} médias, enquanto as ligações \ce{C=O} do
dióxido de carbono são mais fortes que a média; os valores medidos são
para a água líquida, cuja condensação libera mais energia que a reação
entre gases; e o método trata cada ligação como independente das
vizinhas. A estimativa continua útil: ela dá o sinal certo (\omterm{def:g11:reaction-energy:exothermic}{exotérmica})
e a ordem de grandeza certa sem nenhuma medida.
\end{solution}

\begin{solution}{exo:g11:reaction-energy:15}
Rompidas: $4 \times 464 = \qty{1856}{kJ}$; formadas: $2 \times 436 + 498 =
\qty{1370}{kJ}$; $E_r \approx \qty{+486}{kJ}$: \omterm{def:g11:reaction-energy:exothermic}{endotérmica}, logo é
preciso fornecer energia. Queimar o hidrogênio depois devolve essa
energia: o hidrogênio armazena energia produzida em outro lugar, ele não
cria energia.
\end{solution}

\begin{solution}{pb:g11:reaction-energy:1}
\textbf{1.} \ce{CH4 + 2O2 -> CO2 + 2H2O}.

\textbf{2.} \ce{C2H6O + 3O2 -> 2CO2 + 3H2O}.

\textbf{3.} \ce{2C8H18 + 25O2 -> 16CO2 + 18H2O}.

\textbf{4.} \ce{2H2 + O2 -> 2H2O}.

\textbf{5.} O hidrogênio.

\textbf{6.} \qty{16.0}{g/mol}, \qty{46.0}{g/mol}, \qty{114.0}{g/mol},
\qty{2.0}{g/mol}.

\textbf{7.} Em \si{kJ} por \si{kg}, e depois \si{MJ/kg}: metano
$890.7 / 0.0160$, \qty{55.7}{MJ/kg}; etanol $1367.6 / 0.0460$,
\qty{29.7}{MJ/kg}; octano $5470 / 0.1140$, \qty{48.0}{MJ/kg}; hidrogênio
$285.8 / 0.0020$, \qty{142.9}{MJ/kg}.

\textbf{8.} Hidrogênio > metano > octano > etanol.

\textbf{9.} O hidrogênio é um gás muito leve: um quilograma dele ocupa
um volume enorme à pressão normal, por isso ele precisa ser comprimido
em tanques pesados de alta pressão.

\textbf{10.} Metano: 4 \ce{C-H}; gás oxigênio: 2 \ce{O=O} rompidas.
Dióxido de carbono: 2 \ce{C=O}; água: $2 \times 2 = 4$ \ce{O-H}
formadas.

\textbf{11.} $4 \times 415 + 2 \times 498 = \qty{2656}{kJ}$.

\textbf{12.} $2 \times 741 + 4 \times 464 = \qty{3338}{kJ}$.

\textbf{13.} $E_r \approx 2656 - 3338 = \qty{-682}{kJ}$, contra
\qty{-890.7}{kJ} medidos: cerca de \qty{23}{\percent} menor em valor
absoluto.

\textbf{14.} As \omterm{def:g11:reaction-energy:bond-energy}{energias de ligação} médias (o \ce{C=O} do dióxido de
carbono é mais forte que a média); a água líquida na medida, gases na
estimativa.

\textbf{15.} $1000 / 890.7 = \qty{1.123}{mol}$.

\textbf{16.} $1.123 \times 44.0 = \qty{49.4}{g}$.

\textbf{17.} Etanol: $1000 / 1367.6 = \qty{0.731}{mol}$, que dão
\qty{1.462}{mol} de dióxido de carbono, \qty{64.3}{g}. Octano:
$1000 / 5470 = \qty{0.183}{mol}$, que dão \qty{1.463}{mol},
\qty{64.4}{g}.

\textbf{18.} O metano: ele tem o maior número de \omterm{def:g7:atoms-and-molecules:atom}{átomos} de hidrogênio
por \omterm{def:g7:atoms-and-molecules:atom}{átomo} de carbono (4 contra 2.25 no octano), e queimar hidrogênio dá
energia sem dióxido de carbono.

\textbf{19.} Do modo como o hidrogênio é produzido: a partir da água,
com eletricidade de fontes renováveis, ele quase não libera dióxido de
carbono; produzido a partir do metano, ele libera dióxido de carbono na
fábrica em vez de no escapamento.

\textbf{20.} O metano libera cerca de \qty{49}{g} de dióxido de carbono
por megajoule.
\end{solution}
